%Homework 1 
%The file will not compile. Your job is to find the errors and fix them so it produces a clean PDF. Bring your working version to class.

\documentclass[12pt]{article}

\usepackage{enumerate}

\title{Intro to \LaTeX{}\\[4pt] \large Intro to \LateX}
\author{Danish Puri\\ }
\date{\today}

\begin{document}

\maketitle

\tableofcontent %There is a missing letter here 
\pagebreak

\section{Math}

\LaTeX{} (% Think why do we need {}) 
is famous for math. You write math inside math mode. This section shows the basics.

\subsection{Math Modes}
% Think what $ sign does 
\subsubsection{Inline Math Mode (\$)}
Inline math sits inside a sentence. You wrap the math with one dollar sign on each side.

The relationship between three sides of a right angle triangle is a^2+b^2=c^2.
%Can we use $$ for the math equation?

\subsubsection{Display Math Mode}
Display math sits on its own line. \LaTeX{} centers it. You wrap the math with two dollar signs on each side.

The relationship between three sides of a right angle triangle is
$$a^2+b^2=c^2
%So, I opned the env \$\$. Is it closed? 

\subsection{Superscript}
The caret symbol \verb|^| makes a superscript. Braces \verb|{ }| group more than one character.
$$x^2$$
$$x^31$$ %Does this complile? if yes, does it produce what we want?
$$3x^{x^2+5x+2}$$
$$7x^2^3^4$$

\subsection{Subscript}
The underscore \verb|_| makes a subscript. Braces group more than one character.
$$a_0$$
$$a_12$$ %Does this complile? if yes, does it produce what we want?
$$a_{0_{1_{2_{3}}}}$$
The command \verb|\ldots| prints clean dots.
$$a_0, a_1, a_2, a_3 \ldots a_{100}$$

\subsection{Symbols}

\subsubsection{Greek Symbols}
Greek letters need math mode. A capital first letter gives the capital form.
\begin{itemize}
    \item \alpha  %Does this complite despite not being in a math ($) env? 
    \item $\beta}
    \item $\gamma$
    \item $\rho$
    \item $\sigma$
    \item $\delta$
    \item $\epsilon$
    \item $\chi$
    \item $\pi$
    \item $\Pi$
    \item $Area = \pi r^2$
\end{itemize}

\subsubsection{Trigonometry Symbols}
\LaTeX{} has commands for trig functions. The commands print upright letters.
$$ y = sin (x) $$ %does laTeX understand sin?
$$y = \sin (x)$$ %Notice \sin and sin are different 
$$y = \cos (x)$$
$$y = \tan \theta$$
Some books write the inverse like this.
$$y = \sin^{-1} \theta$$
The command \verb|\arcsin| is clearer. Use it instead.
$$y = \arcsin \theta$$

\subsection{Algebraic Symbols}
\begin{itemize}
    \item $\times$
    \item $\div$
    \item \frac{2}{3}
    \item $\frac{\sin \theta}{\cos \theta} = \tan \theta$
    \item $\frac{3x+5}{7x^2+6x+4}$
    \item $\sqrt{5}$
    \item $\sqrt[3]{9}$
    \item $\sqrt[5]{x^2+y^2}$
    \item $\sqrt{5 +\sqrt{x^2+y^2}}$
    \item $\sqrt{\sqrt{\sqrt{\sqrt{\sqrt{\sqrt{2}}}}}}$
\end{itemize}

\subsubsection{Other Symbols}
These symbols appear a lot in calculus and statistics.
$$\int \qquad \oint \qquad \sum \qquad \prod$$

\section{Lists}
\LaTeX{} has two common list types. The \verb|itemize| environment makes bullet lists. The \verb|enumerate| environment makes numbered lists. Me & you can use both.

\subsection{Bullet Lists with Custom Labels}
You can replace a bullet. Put your own label inside square brackets.
\begin{itemize}
    \item[1] Argentina
    \item[2] Switzerland
    \item[3] Morocco
    \item[4] France
    \item Spain
    \item Belgium
    \item Portugal
    \item England
\end{enumerate}

\subsection{Numbered Lists}
The \verb|enumerate| package allows Roman numerals. You can also nest a list inside a list.
\begin{enumerate}[I.] %What package is requried? 
    \item Argentinas
        \begin{enumerate}
            \item Messi
            \item Martinez
            \item Lopez
        \end{enumerate}
    \item Switzerland
    \item Morocco
    \item France
    \item Spain
    \item Belgium
    \item Portugal
    \item England
\end{enumerate}

\subsection{More Custom Labels}
Empty brackets hide the label. Math symbols also work as labels.
\begin{enumerate}
    \item[-] Argentina
        \begin{enumerate}
            \item Messi
            \item Martinez
            \item Lopez
        \end{enumerate}
    \item[-] Switzerland
    \item[*] Morocco
    \item[*] France
    \item[] Spain
    \item[$\to$] Belgium
    \item[=] Portugal
    \item[] England
\end{enumerate}

\pagebreak

\section{Tables}
The \verb|tabular| environment builds a table. The letters \verb|l|, \verb|c| and \verb|r| set the column alignment. The bar \verb!|! draws a vertical line. The symbol \verb|&| separates cells. The command \verb|\\| ends a row. The command \verb|\hline| draws a horizontal line.

\subsection{A Table with a Caption}
The \verb|table| environment wraps the tabular. It adds a caption and a label. The option \verb|[h]| asks \LaTeX{} to place the table here.

\begin{table}[h]
    \centering
    \begin{tabular}{|l|l|}
    \hline
    Country & Player & Goals \\
    \hline
    Argentina & Lionel Messi  & 8 \\
    Norway    & Erling Haaland & 7 \\
    France    & Kylian Mbappe & 7 
    England   & Harry Kane    & 7 \\
    \hline
    \end{tabular}
    \caption{Top scorers in the FIFA World Cup 2026 (sample data)}
    \label{tab:topscorers}
\end{table}

The command \verb|\ref| prints the table number. Table \ref{tab:topscorers} shows sample data.

\subsection{A Plain Tabular}
A tabular without the table wrapper has no caption and no number. A double \verb|\hline| draws a double line.

\begin{center}
\begin{tabular}{|l|l|l|}
\hline
Country & Player & Goals \\
\hline \hline
Argentina & Lionel Messi  & 8 \\
Norway    & Erling Haaland & 7 \\
France    & Kylian Mbappe & 7 \\
England   & Harry Kane    & 7 \\
\hline
\end{tabular}
\end{center}

\section{Text Formatting}

\subsection{Bold}
This is a textbf{bold} text.

\subsection{Italics}
This is an \textit{italicized} text.

\subsection{Nested Styles}
This text is \textbf{\textit{bold and italicized}}.

\subsection{Underline}
This text is \underline{underlined.

\subsection{Font Sizes}
\LaTeX{} has size commands. They run from \verb|\tiny| to \verb|\Huge|.

This is a \begin{huge}huge\end{huge} word.

This is a \begin{Huge}Huge\end{Huge} word.

\subsection{Links}
The \verb|hyperref| package makes clickable links. The command \verb|\href| takes two parts. The first part is the address. The second part is the visible text.

Today's notes live at \href{https://github.com/danish-puri/LaTeX}{this GitHub page}.

The command \verb|\url| shows the raw address instead. See \url{https://github.com/danish-puri/LaTeX}.

\section{Graphics}
Images need the \verb|graphicx| package. This guide already loads it in the preamble.

\subsection{Step 1: Add the Image File}
Upload your image to your project folder. Overleaf has an upload button on the top left. Keep the image next to your \verb|.tex| file. Let the image be \verb|img.jpg|.

\subsection{Step 2: Insert the Image}
The command below inserts the image. Remove the percent sign after you upload \verb|img.jpg|.

\includegraphics[width=0.5\textwidth]{img.jpg}

The option \verb|width=0.5\textwidth| makes the image half as wide as the text. A fixed size like \verb|width=5cm| also works. The option \verb|scale=0.5| shrinks the image to half size. The option \verb|angle=90| rotates it.

\subsection{Step 3: Use the Figure Environment}
The \verb|figure| environment adds a caption and a number. It also lets you refer to the image later. Remove the percent signs after you upload the image.

\begin{figure}[h]
     \centering
     \includegraphics[width=0.5\textwidth]{img.jpg}
     \caption{My first image in LaTeX}
     \label{fig:first}
\end{figure}

The command \verb|\ref{fig:first}| prints the figure number inside your text.

\section{Sectioning and the Table of Contents}
\LaTeX{} builds structure with sectioning commands. Each command makes a heading at a different level.

\begin{itemize}
    \item \verb|\section| makes a main heading.
    \item \verb|\subsection| makes a heading one level down.
    \item \verb|\subsubsection| goes one more level down.
\end{itemize}

The command \verb|\tableofcontents| collects every heading. It prints them at the start with page numbers. Compile the document twice. The first run collects the headings. The second run prints the table.

The next two sections show this structure in action. Find them in the table of contents at the start of this guide.

%You can ignore the rest of doocument. We did not cover symbols for limits, integation.. in class. Pherhaps, watch world up :) 

\section{Polynomials}
A polynomial is a sum of powers of $x$. The degree is the highest power.

\subsection{Linear}
A linear polynomial has degree one. Its graph is a straight line.
$$y = mx + c$$
Here $m$ is the slope and $c$ is the intercept.

\subsection{Quadratic}
A quadratic polynomial has degree two. Its graph is a parabola.
$$y = ax^2 + bx + c$$
The quadratic formula gives the roots.
$$x = \frac{-b \pm \sqrt{b^2 - 4ac}}{2a}$$

\subsection{Cubic}
A cubic polynomial has degree three.
$$y = ax^3 + bx^2 + cx + d$$

\section{Calculus}

\subsection{Limits}
The command \verb|\lim| prints the limit symbol. The arrow uses \verb|\to|.
$$\lim_{x \to 0} \frac{\sin x}{x} = 1$$

\subsection{Derivative}
The derivative measures the rate of change.
$$f'(x) = \lim_{h \to 0} \frac{f(x+h) - f(x)}{h}$$
The power rule is the most common rule.
\frac{d}{dx} x^n = n x^{n-1}

\subsection{Integration}
The command \verb|\int| prints the integral sign. The limits use subscript and superscript.
$$\int x^2 \, dx = \frac{x^3}{3} + C$$
$$\int_0^1 x^2 \, dx = \frac{1}{3}$$

\subsection{Differential Equations}
A differential equation contains a function and its derivatives.

\subsubsection{First Order Differential Equations}
A first order equation contains the first derivative only.
$$\frac{dy}{dx} + P(x)\,y = Q(x)$$

\subsubsection{Second Order Differential Equations}
A second order equation contains the second derivative.
$$a\,\frac{d^2y}{dx^2} + b\,\frac{dy}{dx} + c\,y = 0$$

\section{Where to Go Next}
Practice is the best teacher. Rebuild your last assignment in \LaTeX{}. Use this guide as your reference.

These free resources will help you go further.
\begin{itemize}
    \item Overleaf documentation: \url{https://www.overleaf.com/learn}
\end{itemize}

\end{docment}
